\begin{qfmethod}[Private investment methods and fund structures: closed-end limited partnerships, open-ended and evergreen funds, funds of funds, separate mandates, co-investment and secondary vehicles, and when each fits]\label{meth:private-fund-structures}
\qffield{Statement} \emph{Closed-end limited partnerships} are the standard \citep[pp.~61--63]{demaria-private-markets-2021}:
\begin{itemize}
\item LPs commit at inception; the GP closes the fund at its target size, calls capital during an investment period of three to five years and distributes proceeds as assets are sold;
\item lives are about eight years for direct lending, ten for private equity and real estate, thirteen for funds of funds and timber, fifteen for infrastructure, often with extensions;
\item there are no redemptions: an LP can leave only by selling its stake on the secondary market, at a negotiated price, typically 15--30 percent below NAV;
\item funds are blind pools: the GP sets the strategy, scope and criteria but not the assets in advance.
\end{itemize}
\emph{Alternatives} (pp.~60--66, 95--98):
\begin{itemize}
\item \emph{Open-ended or evergreen} funds have no end date and can be listed. They suit long-lived, income-producing assets and investors who want continuous exposure, but must manage redemptions against illiquid holdings.
\item \emph{Funds of funds} diversify across managers and vintages, give access for modest sums and select managers, at a second layer of fees and carry. Demaria reports few losses across 349 private-equity funds of funds.
\item \emph{Separate mandates} (segregated accounts) are tailored to one investor, discretionary or advisory, from about EUR 30 million, and reduce fees and the number of GP relationships.
\item \emph{Co-investment} vehicles replicate a fund's deals at lower cost for its large investors.
\item \emph{Secondary} funds buy existing LP stakes, often at discounts, shortening the J-curve; Harvard faced discounts of about 50 percent selling during the 2008--2009 crisis \citep[p.~595]{s7-ang}.
\item \emph{Listed vehicles} trade daily but carry the volatility of their share price and premiums or discounts to NAV \citep[p.~596]{s7-ang}.
\end{itemize}

\qffield{Practice} \emph{Matching structures to investors} (original):
\begin{itemize}
\item a \$500 million family office starting private equity: a fund of funds or a diversified mandate across three vintages, then direct funds once it has relationships;
\item a large pension with an in-house team: primary funds with co-investment rights and a separate mandate for small-market buyouts;
\item an insurer needing income and limited capital charges: direct lending funds or a mandate;
\item an investor overcommitted after a market fall: sell older fund stakes on the secondary market, accepting the discount, rather than default on calls.
\end{itemize}

\qffield{Conditions} Lives, discounts and minimum sizes are those Demaria reports; they vary by market and period.

\qffield{Pitfalls} Treating an evergreen fund's NAV as a price at which all investors can redeem at once. Paying fund-of-funds fees for diversification the investor could reach alone.

\qffield{See also} \hyperref[def:alternative-investment-categories]{methods of investing}; \hyperref[meth:alternatives-allocation-approaches]{liquidity planning}.

\end{qfmethod}
